\documentstyle[%


righttag%


,11pt%


,amssymb%


,russian%               remove in Eng.


,izvru%                 Set izven% in Eng.


]{amsart}





%





\newcommand{\tts}[1]{}





\theoremstyle{plain}


\theorembodyfont{\it}


\newtheorem{thmnonum}{\hskip\parindent Теорема}


\renewcommand{\thethmnonum}{}





\theoremstyle{definition}


\theorembodyfont{\rm}


\newtheorem{notenonum}{\hskip\parindent Замечание}


\renewcommand{\thenotenonum}{}





%\hoffset=-15mm%                  For Eng.


\setcounter{page}{57}


\god{2005}


\nomer{9~(520)} %                        Remove in Eng.


%\nomer{Vol.\,49, No.\,8}%               Set in Eng.


%\str{\pageref{begin}--\pageref{end}}%  Set in Eng.


\udk{517.988}





\author{\it В.Я.\,Прудников}


% NB:   ^^ remove in Eng.


\title{О коэрцитивности выпуклых функционалов}





\date{}


%\pagestyle{myheadings}%         Set in Eng.


\pagestyle{plain} %              Remove in Eng.


\begin{document}


%\thispagestyle{plain}%          Set in Eng.


\maketitle


\label{begin}








Для существования решения вариационных неравенств на соответствующий 


функционал часто налагается требование его коэрцитивности [1], [2].


Если функционал не коэрцитивен, то существует, как отмечено,


например, в [3], конструктивный метод исследования вариационных 


неравенств, основанный на построении коэрцитивных задач. 





В данной работе дан критерий коэрцитивности выпуклых по направлению 


функционалов на замкнутом выпуклом конусе.





Будем предполагать, что $K$ --- замкнутый выпуклый конус из 


рефлексивного банахова пространства $X$, $S=\{u \in X: \|u\|=1\}$, 


$B=\{u \in X: \|u\|<1\}$, $SK=S \cap K$, $BK = B \cap K$.


Под коэрцитивностью функционала $p: K \to R$


понимается существование предела


$p(u) \to +\infty$, если $\|u\| \to +\infty$, $u\in K$. 








\begin {thmnonum}


Пусть функционал $p: K \to R$ таков, что


\begin{itemize}


\item[1)] функция $t \to p(at)$ выпукла в интервале $(0,+\infty)$ 


для любого $a \in SK;$


\item[2)] функционал $a \to p(at)$ слабо полунепрерывен снизу на


$\overline {BK}$ для любого $t\in(0,+\infty)$.


\end{itemize}


Тогда


$$


(p\text{ коэрцитивен на $K$ }) \Longleftrightarrow


\bigg(\inf_{a \in SK}


\lim_{t \to +\infty} \frac {p(at)}{t} > 0 \bigg).


$$


\end{thmnonum}





Для доказательства теоремы докажем три леммы, предполагая выполнимость 


условий теоремы.





\begin {lem}


$(p$ коэрцитивен на $K)$ $\Longleftrightarrow


\big(\lim\limits_{t \to +\infty }\inf\limits_{a \in SK}


\frac {p(at)}{t} > 0 \big)$. 


\end {lem}





\begin{pf} {\bf Необходимость.}


Из существования предела


$\lim\limits_{\|u\| \to +\infty} p(u)= +\infty$


для данного $A>0$ следует существование такого числа $R_A >0$, что при


$\|u\| > R_A$ выполнимо неравенство $p(u) > A$. Полагая здесь $u=ta$,


где $t=\|u\|$, $a \in SK$, получим неравенство $p(ta) > A$,


справедливое при любых $t>R_A$ равномерно относительно $a \in SK$,


а потому


$$


\inf_{a \in SK} p(ta) > A, \quad t> R_A,


$$


что и означает существование предела


$\lim\limits_{t \to +\infty} \inf\limits_{a \in SK} p(at) = +\infty$.


Поэтому существует такое $t_0>0$, что


$\inf\limits_{a \in SK} p(t_0a) > 1 + p(\theta )$.


Функция $t \to p(ta)$ выпукла по условию теоремы, значит, функция


$t \to \frac {p(at)-p(\theta )}{t}$ не убывает в $(0,+\infty)$.


Но тогда


$$


\frac {p(at)-p(\theta)} {t} \ge \frac {p(at_0)-p(\theta )}{ t_0},


\quad t > t_0,


$$


для любых $a \in SK$, а потому


$$


\inf_{a \in SK}\frac{ p(at)-p(\theta )}{t} \ge


\inf_{a \in SK}\frac {p(at_0)-p(\theta )}{ t_0} \ge


\frac {1}{t_0}, \quad t > t_0,


$$


и


$\lim\limits_{t \to +\infty}


\inf\limits_{a \in SK} \frac {p(at)-p(\theta )}{ t}> 0$.


Учитывая, что $\lim\limits_{t \to +\infty}\frac {p(\theta)}{t}=0$,


получим неравенство


$\lim\limits_{t \to +\infty}\inf\limits_{a \in SK}


\frac{p(at)}{t}>0$.


\smallskip





{\bf Достаточность.} Пусть


$\lim\limits_{t \to +\infty}\inf\limits_{a \in SK}\frac {p(at)}{t}>0$.


Тогда существуют числа $\beta>0$, $t_0 > 0$ такие,


что $\inf_{a \in SK} p(at) > t\beta$ при $t>t_0$.


Следовательно, для всех $a \in SK$ независимо от $t>t_0$


справедливо неравенство $p(at) > t\beta$.


Для элемента $u \in K$, $t = \|u\|$, $t>t_0$, существует $a \in SK$


такой, что $u=\|u\|a$, а потому $p(u)>\|u\|\beta$, $\|u\|>t_0$. 


Отсюда $\lim\limits_{\|u\| \to +\infty} p(u)=+\infty$.


\end{pf} 





\begin {lem}


$(p$ коэрцитивен на $K)$ $\Longleftrightarrow$ $($существует


$t_0>0 :\inf\limits_{a \in SK} p(at_0) > p(\theta))$.


\end{lem}





\begin{pf} {\bf Необходимость} следует из леммы 1.





{\bf Достаточность.} Имеем


$$


\lim_{t \to +\infty}\inf_{a \in SK}\frac {p(at)} {t} =


\lim_{t \to +\infty}\inf_{a \in SK}\frac {p(at)-p(\theta)}{ t} \ge


\inf_{a \in SK}\frac {p(at_0)-p(\theta )}{ t_0}> 0.


$$


Согласно лемме 1 функционал $p$ коэрцитивен на $K$.


\end{pf}





\begin {lem}


Если


$\inf\limits_{a \in SK} \lim\limits_{t \to +\infty} \frac{p(at)}{t} \ge 0$,


то $\forall \varepsilon > 0$ функционал


$p(u) + \varepsilon \|u\|^\alpha$, $\alpha \ge 1$ коэрцитивен на $K$.


\end{lem}





\begin{pf}


Без ограничения общности рассуждений будем


считать, что $p(\theta) =0$. Для $n \in N$ существует элемент $a_n$ такой,


что $\|a_n\|=1$, $a_n \in K$ и


$$


\inf_{a \in SK} \frac{p(an)+\varepsilon n^\alpha}{n} >


\frac {p(a_n n)}{n} + \varepsilon n^{\alpha -1} -\frac {1}{n}.


$$


Функция $t \to\frac{p(a_n t)}{t}$ не убывает в $(0,+\infty)$, 


поэтому при $n \ge t $ выполнено неравенство


$$


\inf_{a \in SK} \frac{p(an)+\varepsilon n^\alpha}{n} >


\frac {p(a_n t)}{t} + \varepsilon t^{\alpha -1} -\frac {1}{n}.


$$


Вследствие рефлексивности пространства $X$ из последовательности


$\{a_n\}$ выделяем слабо сходящуюся к некоторому элементу


$a_0 \in \overline {B}K$


подпоследовательность, за которой сохраняем обозначение $\{a_n\}$.


Поскольку функционал $p(at)$ слабо полунепрерывен снизу на


$ \overline {B}K$, получим неравенство


$$


\lim_{n \to +\infty}


\inf_{a \in SK} \frac{p(an)+\varepsilon n^\alpha}{n}\ge


\frac {p(a_0 t)}{t} + \varepsilon t^{\alpha -1}.


$$


Если $a_0=\theta$, то из этого неравенства следует


$$


\lim_{n \to +\infty}


\inf_{a \in SK} \frac{p(an)+\varepsilon n^\alpha}{n}> 0.


$$


Пусть $a_0 \ne \theta$, тогда


$$


\lim_{n \to +\infty}\inf_{a \in SK} \frac{p(an)+\varepsilon n^\alpha}{n}


\ge\lim_{t \to +\infty}\frac{p(a_0 t)}{t} +\varepsilon =


%


\|a_0\| \lim_{t \to +\infty}


\frac {p\big(\frac{a_0}{\|a_0\|} t \|a_0\| \big)}


{t \|a_0\|}+ \varepsilon \ge \|a_0\|


\inf_{a \in SK}\lim_{t \to +\infty}


\frac{p(at)}{t} +\varepsilon > 0.


$$


Таким образом, существует такое $n_0 \in N$, что


$\inf\limits_{a \in SK}(p(an_0)+\varepsilon n_0^\alpha) > 0$.


Согласно лемме 2 функционал $p(u) + \varepsilon \|u\|^\alpha$


коэрцитивен на $K$.


\end{pf}








\begin{pf*}{Доказательство теоремы}


{\bf Необходимость.} Пусть $p$ коэрцитивен на $K$. Тогда из леммы~1 


следует $\lim\limits_{t \to +\infty}\inf\limits_{a \in SK}


\frac{p(at)}{t} >0$. Существует $\varepsilon > 0$ такое, что


$\lim\limits_{t \to +\infty}\inf\limits_{a \in SK}


\frac{p(at)}{t} > \varepsilon$. Таким образом, для любого


$a \in SK$ \ $\lim\limits_{t \to +\infty}\frac{p(at)}{t}


\ge \varepsilon $,


а потому $\inf\limits_{a \in SK}\lim\limits_{t \to +\infty}


\frac{p(at)}{t} \ge \varepsilon > 0$.


\smallskip





{\bf Достаточность.} Предполагаем без ограничения общности


рассуждений, что $p(\theta)=0$. Существует число $\varepsilon > 0$


такое, что


\begin{equation}


\inf_{a \in SK}\lim_{t \to +\infty}\frac{p(at)-\varepsilon t}{t} \ge 0.


\end{equation}


Функционал $p(u)-\varepsilon \|u\|$ таков, что функция


$t \to p(at)-\varepsilon t$ выпукла для всякого $a \in SK$ в $(0,+\infty)$,


а функционал $a \to p(at)-\varepsilon t$ слабо полунепрерывен снизу на


$\overline{B}K$. В силу (1) и согласно лемме 3 функционал


$p(u)-\varepsilon \|u\| + \varepsilon \|u\|$ коэрцитивен на $K$.


\end{pf*}





\begin{notenonum}


Леммы 1, 2 справедливы без требований


рефлексивности пространства $X$ и слабой полунепрерывности функционала


$a \to p(at)$, $a \in \overline{B}K$.


\end{notenonum}








\section*{Примеры}





1. {\it Задача Синьорини} (напр., [4]).


Пусть $K=\{ u \in H^1(E) : u \geq 0$ почти всюду на $\Gamma\}$, где


$\Gamma$ --- достаточно гладкая граница ограниченной области 


$E \subset R^n$, $\mathrm{mes}\,E=1$. Обозначим


$$


b(u,u)=\int_E|\nabla u|^2dE, \quad


(u,f) =\int_EufdE,\quad


p(u)=\frac{1}{2} b(u,u) -(f,u), \quad


\|u\|=\sqrt{b(u,u)}+\sqrt{(u,u)}.


$$


Имеем


$$


\inf\bigg\{ \lim_{t \to +\infty}\frac{p(ta)}{t}:


b(a,a) \ne 0, \ a \in SK\bigg\}= +\infty.


$$


Пусть $b(a,a)=0$, $a \in SK$, тогда $a=1$. Значит,


$\lim\limits_{t \to +\infty}\frac{p(t)}{t}=-\int\limits_EfdE$.


Таким образом, условие $\int\limits_EfdE < 0$


необходимо и достаточно для коэрцитивности функционала $p$.


\smallskip





2. {\it Модельная задача с трением} [4].


Пусть $K=H^1$, $p(u)=\frac12 b(u,u) + j(u) - (f,u)$, где


$j(u)=\int\limits_\Gamma g|u|d\Gamma$, $g \ge 0$ на $\Gamma$.


Для $a \in S$, $t > 0$ имеем


$\frac{p(ta)}{t} = \frac {t}{2} b(a,a) + j(a) - (f,a)$.


Если $b(a,a)>0$, то


$\lim\limits_{t \to +\infty}\frac{p(ta)}{t} = +\infty$.


Пусть $b(a,a)=0$, $a \in S$. Тогда $a= \pm 1$, а потому


$\frac{p(ta)}{t} = j(1) - \pm (f,1)$.


Следовательно, неравенство


$\int\limits_\Gamma gd\Gamma >\Big|\int\limits_EfdE\Big|$


необходимо и достаточно для коэрцитивности функционала~$p$.


Заметим, что в [4] это условие приводится лишь как достаточное.





\begin{thebibliography}{9}





\bibitem{1}


Экланд Н., Темам Р. {\it Выпуклый анализ и вариационные проблемы}.


-- М.: Мир, 1979. -- 400~с.


\bibitem{2}


Куфнер А., Фучик С. {\it Нелинейные дифференциальные уравнения}.


-- М.: Наука, 1988. -- 304~с.


\bibitem{3}


Лапин А.В. {\it Решение вариационных неравенств с нелинейными 


полукоэрцитивными операторами} // Вычисл. процессы и системы. -- 


М.: Наука, 1986. -- T.\,4. -- С.\,219--264.


\bibitem{4}


Намм Р.В. {\it Введение в теорию и методы решения вариационных


неравенств}. -- Хабаровск, Изд-во ХГТУ, 1999. -- 72~с.





\end{thebibliography}





\vskip 2cm





\post{\it Хабаровский государственный}{\it Поступила}


\post{\it технический университет}{03.06.2003}





\label{end}


\end{document}


